\chapter{Synthetic Document Generation: French Clinical Application}
\label{chap:french-application}

The previous chapters (\autoref{chap:synthetic} and \autoref{chap:privacy}) validated our synthetic generation methodology on English clinical data from \gls{mimic}-III.
A natural question follows: does this approach generalize to other languages and healthcare systems?

This chapter presents an industrial application of our methodology to French clinical documents, conducted in collaboration with Assistance Publique -- Hôpitaux de Paris (\gls{aphp}).
This work was led by Riccardo Tripodi as part of his master's thesis at Politecnico di Milano, adapting our generation pipeline to the French clinical context.
The results are encouraging and pave the way for larger-scale deployment within the \gls{aphp} infrastructure.

\minitoc

\section{Motivation and Context}

While the experiments presented in previous chapters demonstrate the effectiveness of our approach on English clinical data from \gls{mimic}-III, a critical question remains: does this methodology generalize to other languages and healthcare systems?
This consideration is particularly relevant given that the majority of clinical \gls{nlp} research focuses on English, leaving non-English healthcare systems underserved despite facing similar---or often greater---data scarcity and privacy challenges.

To address this question and validate our approach in a more applied setting, we conducted experiments on French medical reports in collaboration with Assistance Publique -- Hôpitaux de Paris (\gls{aphp}).
Beyond the language shift, this application involves clinical documents from a different medical context: while \gls{mimic}-III contains intensive care unit discharge summaries from a US academic hospital, the French corpus comprises diverse medical reports reflecting routine clinical practice across multiple hospital departments.

From a practical perspective, \gls{aphp} faces concrete challenges that motivate this work.
The hospital system seeks to deploy small, efficient models capable of classifying clinical documents for tasks such as automated \gls{icd} coding.
However, annotated training data remains sparse, and privacy regulations prevent direct use of patient records for model development.
Synthetic document generation addresses both constraints simultaneously.

Beyond immediate operational needs, this collaboration pursues broader research objectives.
First, generating high-quality synthetic documents enables sharing clinical data with the research community without compromising patient privacy---a significant step toward reproducible medical \gls{nlp} research in French.
Second, this work lays the groundwork for a longer-term goal: generating entirely synthetic patient records that preserve the statistical and clinical properties of real data while containing no actual patient information.

This setting presents several distinctive characteristics:

\begin{itemize}
    \item \textbf{Language and clinical conventions}: French clinical writing follows distinct conventions, with specific medical terminology, abbreviation patterns, and narrative structures that differ from English documentation practices.

    \item \textbf{Coding system}: French hospitals use \gls{icd}-10 codes rather than the \gls{icd}-9 codes present in \gls{mimic}-III, representing a different classification granularity.

    \item \textbf{Regulatory context}: European healthcare data is governed by \gls{gdpr}, imposing stringent requirements on data processing and transfer that make privacy-preserving approaches particularly valuable.

    \item \textbf{Fully synthetic seed corpus}: Unlike our \gls{mimic}-III pipeline, which required a small set of real pseudo-anonymized documents to bootstrap the generator (\autoref{chap:synthetic}), the French pipeline eliminates this dependency entirely.
    The seed corpus is generated synthetically from curated mappings between \gls{icd}-10 codes and medical keywords, using a large medical language model, without any real clinical document involved.
\end{itemize}

This cross-lingual validation thus serves multiple purposes: testing whether our methodology transfers across languages and clinical contexts, providing a practical solution for data-scarce healthcare systems, and opening pathways toward shareable synthetic medical corpora.

\section{Experimental Setup}

\paragraph{Dataset.}
The reference corpus used for scoring during \gls{dpo} refinement consists of fictive clinical reports written by practicing physicians across multiple hospital departments and medical specialties.
These documents are clinically realistic but contain no actual patient information, ensuring full reproducibility while eliminating any risk of privacy leakage.
Evaluation is performed on a strictly held-out test set of real clinical reports annotated by experts, fully disjoint from all data used in the pipeline.

\paragraph{Seed corpus generation.}
Prompts are constructed from a curated database linking \gls{icd}-10 codes to medically grounded French keywords describing symptoms, anatomy, treatments, and procedures.
Codes are sampled according to empirical hospital frequency distributions, and a constrained subset of associated keywords is selected for each prompt.
MedGemma-27B-IT~\citep{sellergren2025medgemma}, a large multilingual medical language model, generates \num{20000} synthetic clinical reports following a fixed eight-section narrative structure: \emph{Motif d'hospitalisation}, \emph{Antécédents}, \emph{Mode de vie}, \emph{Histoire de la maladie}, \emph{Examen clinique}, \emph{Examens complémentaires}, \emph{Évolution pendant l'hospitalisation}, and \emph{Conclusion}.
This seed corpus is entirely independent of real patient records, in contrast to the pseudo-anonymized seed documents used in our \gls{mimic}-III experiments (\autoref{chap:synthetic}).

\paragraph{Generator model.}
We employ Qwen3-4B-Instruct~\citep{yang2025qwen3} as the base generator model, selected for its strong multilingual capabilities and competitive performance on French text generation tasks.
Compared to Mistral-7B used in our \gls{mimic}-III experiments, Qwen3-4B offers a favorable trade-off between model size and generation quality for non-English languages.
The model is fine-tuned using \gls{lora}~\citep{hu2022lora} with rank $r=32$ and scaling factor $\alpha=64$, applied to all attention projection layers.
Training follows the same two-stage approach: initial supervised fine-tuning on the \num{20000} synthetic seed documents, followed by three iterations of \gls{dpo} alignment.

\paragraph{Privacy-bounded refinement.}
During \gls{dpo} iterations, the generator produces candidate reports that are scored against the fictive reference documents inside the hospital infrastructure.
For each keyword prompt, $N=4$ candidates are generated using a mixed temperature strategy: one sample at $\tau = 0.2$ to favor coherent outputs, and three samples with temperatures drawn from $\mathcal{U}(0.6, 0.9)$ to encourage diversity.
Keywords are extracted from reference reports using QuickUMLS~\citep{soldainiQuickUMLSFastUnsupervised2016}.
The complete scoring mechanism is detailed in \autoref{sec:french-scoring}.

\paragraph{Downstream task.}
We evaluate synthetic document quality through \gls{icd}-10~\citep{who_icd10} multilabel classification, mirroring our \gls{icd}-9 experiments on \gls{mimic}-III.
A ModernBERT-large classifier~\citep{warner2025smarter} is trained exclusively on synthetic datasets of varying sizes (\num{10000}, \num{20000}, \num{50000} documents) and evaluated on the held-out real test set.
Classification targets are defined by the top-$k$ most frequent \gls{icd}-10 codes with $k \in \{20, 50, 100\}$.
At the \num{10000} scale, this yields approximately \num{350} samples per code for top-20 but only \num{150} for top-100.
The average number of codes per document increases with $k$: $1.52 \pm 0.82$ for top-20, rising to $2.10 \pm 1.14$ for top-100.
Generated reports average approximately $800 \pm 240$ words across refinement stages.

\section{Scoring Mechanism}
\label{sec:french-scoring}

While our \gls{mimic}-III experiments relied primarily on SemScore for document quality assessment, the French application employs a composite scoring mechanism combining two complementary components.

\paragraph{Semantic similarity.}
The first component measures cosine similarity between synthetic and reference documents using BioLORD-2023-M~\citep{remy2024biolord} embeddings.
Unlike general-purpose embedding models, BioLORD integrates clinical knowledge graph information, making it particularly suitable for assessing semantic fidelity in medical documents across languages.

\paragraph{\gls{llm}-as-judge.}
The second component uses Qwen3-30B-A3B-Instruct~\citep{yang2025qwen3} as an \gls{llm}-as-judge to evaluate medical correctness and content alignment of the generated candidate with respect to the reference report.

\paragraph{Composite score.}
The final quality score combines these two components through simple averaging:
\begin{equation}
    S_{composite} = \frac{1}{2}\left(S_{similarity} + S_{judge}\right)
\end{equation}
This composite score guides candidate selection during \gls{dpo} alignment, with higher-scoring documents serving as preferred examples and lower-scoring ones as rejected examples.
Compared to the SemScore-only approach used on \gls{mimic}-III, this two-faceted evaluation provides a stronger training signal by capturing both semantic fidelity and clinical plausibility.

\section{Example of Generated Report}

\autoref{fig:french-example} presents an excerpt of a synthetic French clinical report generated at the end of the pipeline (after DPO3).
The document follows the eight-section narrative structure and demonstrates clinically coherent content with appropriate medical terminology.

\begin{figure}[htb]
\centering
\begin{tcolorbox}[colback=gray!5,colframe=gray!50,boxrule=0.5pt,arc=2pt,left=8pt,right=8pt,top=6pt,bottom=6pt,fontupper=\small]
\textbf{Motif d'hospitalisation:}
Mme Dupont, 68 ans, est admise en hospitalisation de jour le 15/03/2024 pour traitement d'une angiodermite nécrotique extensive surinfectée. Cette lésion cutanée chronique, initialement diagnostiquée il y a plusieurs années, s'est aggravée de manière significative, nécessitant une prise en charge intensive et une surveillance rapprochée.\\[0.5em]

\textbf{Histoire de la maladie:}
L'angiodermite nécrotique initiale de Mme Dupont a connu des épisodes récurrents et des aggravations progressives au fil des ans. Depuis environ 6 mois, une lésion cutanée extensive, initialement localisée au niveau des membres inférieurs, s'est étendue progressivement et s'est caractérisée par une nécrose cutanée et une infection bactérienne sévère. La lésion est devenue très douloureuse, avec des signes d'inflammation marquée (érythème, chaleur, \oe{}dème) et une ulcération profonde. La patiente a consulté son dermatologue qui a initié un traitement antibiotique (Ceftriaxone) et une antibiothérapie topique, mais l'aggravation persiste.\\[0.5em]

\textit{[\ldots]}\\[0.5em]

\textbf{Mots-clés:} Angiodermite nécrotique extensive surinfectée, personne vivant seule à son domicile.
\end{tcolorbox}
\caption{Excerpt of a synthetic French clinical report generated after three \gls{dpo} iterations. The document follows the standardized eight-section structure and demonstrates clinically coherent content.}
\label{fig:french-example}
\end{figure}

\section{Results}
\label{sec:french-results}

Table~\ref{tab:french-results} presents classification performance across dataset sizes and \gls{dpo} iterations.
We report macro F1 scores on the held-out test set, with relative gains computed against the Step0 baseline (\gls{sft} only, no \gls{dpo} alignment).

\begin{table}[htb]
\centering
\begin{tabular}{@{}llccccr@{}}
\toprule
\textbf{Dataset} & \textbf{Top-k} & \textbf{Step0} & \textbf{DPO1} & \textbf{DPO2} & \textbf{DPO3} & \textbf{Gain} \\
\midrule
\multirow{3}{*}{10K}
    & 20  & 0.389 & \textbf{0.422} & 0.422 & 0.418 & +7.4\% \\
    & 50  & 0.250 & 0.294 & 0.302 & \textbf{0.317} & +26.9\% \\
    & 100 & 0.175 & 0.197 & 0.202 & \textbf{0.238} & +36.5\% \\
\midrule
\multirow{3}{*}{20K}
    & 20  & 0.438 & \textbf{0.457} & 0.448 & 0.447 & +2.1\% \\
    & 50  & 0.299 & 0.326 & 0.334 & \textbf{0.347} & +16.2\% \\
    & 100 & 0.195 & 0.230 & 0.210 & \textbf{0.285} & +45.8\% \\
\midrule
\multirow{3}{*}{50K}
    & 20  & 0.445 & \textbf{0.461} & 0.461 & 0.455 & +2.3\% \\
    & 50  & 0.326 & 0.351 & \textbf{0.385} & 0.383 & +17.6\% \\
    & 100 & 0.230 & 0.268 & 0.300 & \textbf{0.316} & +37.2\% \\
\bottomrule
\end{tabular}
\caption{Classification F1 scores on French \gls{icd}-10 multilabel task across dataset sizes and \gls{dpo} iterations.
Relative gains computed against Step0 baseline.}
\label{tab:french-results}
\end{table}

The results reveal consistent patterns that mirror our \gls{mimic}-III findings while demonstrating even stronger improvements from \gls{dpo} alignment.
Relative gains range from +2\% for top-20 to +46\% for top-100, substantially higher than the +3\% to +6\% gains observed on English data.
This difference likely stems from the composite scoring mechanism combining semantic similarity and \gls{llm}-as-judge, which provides a richer training signal for preference optimization compared to SemScore alone.

The relationship between task complexity and \gls{dpo} benefit follows the same pattern as \gls{mimic}-III.
For top-20 classification, a single \gls{dpo} iteration often suffices, with subsequent iterations providing marginal or no additional improvement.
In contrast, top-50 and top-100 configurations benefit from continued alignment through DPO2 and DPO3, confirming that rare label prediction requires higher-quality synthetic documents regardless of language.

Dataset scaling effects remain consistent: 50K models outperform 20K and 10K counterparts across all configurations.
However, the most striking observation is that even the smallest 10K dataset with DPO3 alignment achieves competitive performance, suggesting that quality-focused generation can partially compensate for limited data quantity.

\section{Cross-lingual Comparison}
\label{sec:crosslingual}

To assess whether our methodology exhibits consistent behavior across languages, we compare key patterns between the English (\gls{mimic}-III) and French (\gls{aphp}) experiments.

The qualitative patterns are remarkably consistent across both languages.
In both settings, \gls{dpo} alignment provides clear benefits over the \gls{sft} baseline, with gains concentrated in higher top-k configurations where rare labels dominate.
The diminishing returns pattern for simpler tasks (top-20) appears in both experiments, with DPO1 often providing the best results before subsequent iterations plateau.
Iterative \gls{dpo} through DPO2 and DPO3 proves beneficial for complex tasks (top-100) in both languages, and dataset scaling consistently improves absolute performance.

The quantitative differences are notable: French experiments show substantially larger relative improvements from \gls{dpo} alignment (+36\% to +46\% for top-100) compared to English (+3\% to +6\%).
We attribute this primarily to the composite scoring mechanism employed in the French pipeline, which combines semantic similarity and \gls{llm}-as-judge evaluation.
This two-component scoring likely provides a stronger training signal for preference optimization compared to the SemScore-only approach used on \gls{mimic}-III.

Additional factors may contribute to these differences: the Qwen3-4B model may respond more favorably to \gls{dpo} than Mistral-7B, and \gls{icd}-10 classification may present different challenges than \gls{icd}-9.
However, the consistent qualitative patterns across languages provide strong evidence that the core methodology---keyword-guided generation with iterative preference optimization---transfers effectively regardless of these implementation details.

\section{Discussion \& Conclusion}

These results provide strong evidence that our reward-based synthetic generation approach generalizes effectively beyond English.
The consistent improvement patterns observed across \gls{mimic}-III and \gls{aphp} experiments---despite differences in language, coding systems, and model architectures---suggest that the underlying methodology captures language-agnostic principles of clinical document structure.

From a practical standpoint, this cross-lingual validation is promising for the medical domain.
Healthcare systems worldwide face similar challenges: sensitive patient data that cannot be shared, sparse annotations, and regulatory constraints on data processing.
Our results demonstrate that synthetic document generation offers a viable pathway to address these challenges regardless of language, opening opportunities for medical \gls{nlp} research in underserved linguistic communities.

Looking forward, the collaboration with \gls{aphp} opens perspectives for extending this work.
One objective is the generation of entirely synthetic patient records, leveraging the multimodal data available in hospital information systems to create coherent document collections that span a patient's clinical trajectory.
Another direction concerns long-context generation: current experiments focus on relatively short documents, and scaling to longer clinical narratives such as complete discharge summaries or multi-document patient histories remains to be explored.

\paragraph{Acknowledgements.}
This work on French clinical data was conducted in collaboration with Assistance Publique -- Hôpitaux de Paris (\gls{aphp}).
We thank Riccardo Tripodi for implementing and running the experiments on French medical reports.
